\section{Practical Implications and Checkpoint Selection}
\label{sec:recipe}

The three studies support practical choices about weighting, diagnostics, and evaluation. They do not identify one configuration that is best across domains and evaluation sets.

\paragraph{Specify the intended weighting and measure the resulting capability profile.}
GT weights tokens equally; DT fixes domain weights while retaining length weighting within each domain; DR also weights responses equally. These are different objectives, even with the same prompt quotas. Likewise, the PG/I16/I64 results favor selecting vocabulary feedback from per-domain capability and measured cost, with retained probability mass used to interpret coverage. The present runs do not provide a matched-cost efficiency ranking.

\paragraph{Interpret update geometry with its history and precision.}
Report the parameter representation, distinguish individual steps from cumulative changes, and compare teacher-conditioned updates at a common state. A zero-current-gradient reference helps assess how much alignment depends on shared history. Evidence of overlap or opposing directions should be followed by per-domain loss or capability measurements before inferring functional interference.

\paragraph{Evaluate checkpoint selection on separate questions.}
To test domain retention as a selection criterion, let $s_d(\theta)$ be a development score, with higher values better, and let $\epsilon_d\geq0$ be a prespecified tolerance relative to initialization. Eligible checkpoints satisfy
\begin{equation}
s_d(\theta)\geq s_d(\theta_0)-\epsilon_d
\qquad\text{for every domain }d.
\label{eq:retention}
\end{equation}
We select the eligible checkpoint with maximum equal-domain mean, retaining the initial student if none improves it by at least one percentage point. We compare this rule with unconstrained maximum utility and the final checkpoint, using a two-point domain tolerance fixed before the additional test results. Appendix~\ref{app:completed_selection} gives the development grids and selection outcomes.


The additional test uses separately frozen questions from four domain sources, with 256 questions per domain and a distinct evaluation protocol (Appendix~\ref{app:independent_domain_test}). Relative rankings can change: PG-DR/SGD exceeds Adam by 1.18 mean points on the original benchmarks, but is 0.07 points lower on this test. Retention and unconstrained maximum utility select the same checkpoint on all ten trajectories: update 500 in nine cases and update 250 for PG-DT/Adam. In the latter case, the selected checkpoint scores 43.14\% on the additional test, versus 43.51\% at update 500. The retention constraint therefore has no demonstrated selection benefit here. Because the development benchmarks had been observed previously, this selection study remains exploratory.
